\chapter{Conclusion and Future Work}
\label{chap:conclusion}

\section{Summary of Contributions}
\label{sec:contributions}

This dissertation has delivered six concrete contributions to the intersection of classical artificial intelligence, process mining, and verifiable computation.

\textbf{1. Thirteen classical AI architectures validated and falsified in WebAssembly.}
The \texttt{wasm4pm} cognition layer implements MYCIN, Hearsay-II, SOAR, CBR, Prolog, STRIPS, GPS, DENDRAL, ELIZA, and four autoinstinct breeds (vision, semantics, neurosis, learning) as \texttt{wasm\_bindgen}-exported modules compiled to both browser and Node.js targets. Each breed is not merely implemented but \emph{falsified}: every claim about its behavior is backed by an unfakeable oracle that would fail if the underlying algorithm were replaced by a stub. This constitutes the first comprehensive WebAssembly port of the classical AI canon with adversarial test coverage.

\textbf{2. An OCEL provability layer (ocel.rs, 632 lines) enforced at the run-breed choke point.}
The module \texttt{src/ocel.rs} defines \texttt{OcelEvent}, \texttt{OcelObject}, and \texttt{OcelLog} conformant to the OCEL~2.0 specification, a \texttt{derive\_ocel} function that converts a breed's \texttt{TraceStep} sequence into a structured event log with logical (not wall-clock) timestamps, thirteen \texttt{BreedLifecycleModel} constants encoding the declared lifecycle of each architecture, \texttt{validate\_ocel\_alignment} implementing a DFA phase checker against those models, and \texttt{check\_temporal\_conformance} enforcing strictly monotone logical step ordering. All of this is wired at \texttt{run\_breed} in \texttt{wasm.rs}: a breed whose execution trace does not conform to its declared lifecycle model is refused, not just flagged. This is the first process-mining-grade conformance gate applied to classical AI execution.

\textbf{3. A BLAKE3 receipt chain whose hash coverage includes the OCEL log.}
Every breed execution produces \texttt{output\_hash = blake3(serde\_json(BreedOutput))}, where \texttt{BreedOutput} contains the \texttt{ocel\_log}. The \texttt{run\_id} is derived as \texttt{blake3(breed\,|\,output\_hash)}, and the \texttt{replay\_pointer} is the first sixteen hex characters of \texttt{output\_hash}. Because the OCEL log is inside the hashed payload, the receipt attests to process behavior, not merely to a result value. A receipt cannot be forged without re-executing the breed; a receipt cannot pass with a non-conforming trace.

\textbf{4. A 655-test falsification suite: 480 Rust and 175 TypeScript, with zero failures, providing Rank-4 oracles for all thirteen breeds.}
The Rust suite spans fifteen test files including adversarial bypass attempts (\texttt{breed\_adversarial.rs}, 58 tests), mathematical property invariants (\texttt{breed\_oracle\_gaps.rs}, 31 tests), unfakeable autoinstinct oracles (\texttt{autoinstinct\_adversarial.rs}, 27 tests), formal invariants for STRIPS, SOAR, and CBR (\texttt{strips\_soar\_cbr\_invariants.rs}, 23 tests), and OCEL lifecycle conformance (\texttt{ocel\_conformance.rs}, 7 tests). Key unfakeable properties include: the Prolog grandparent derivation, which requires shared-variable Robinson unification across clauses; MYCIN's continuous CF boundary at 0.2, which a threshold stub cannot fake; Hearsay's opportunistic KSAR ordering, which requires a priority queue sensitive to \texttt{rating = certainty * trigger\_cf}; and SOAR antisymmetry under circular worse-than preferences. The TypeScript suite covers the contract layer, receipt emission, and OCEL alignment across all eleven packages.

\textbf{5. Algorithm fidelity upgrades restoring the substance of the originals.}
Prior WebAssembly ports of classical AI architectures have uniformly simplified. This work reversed that trend. Hearsay received the KSAR priority queue with \texttt{f32::total\_cmp} ordering and stale-KSAR re-rating on blackboard update. SOAR received the full preference algebra (require, prohibit, better, worse, best, indifferent) and bounded subgoal resolution on tie impasse with a depth cap. CBR received the complete 4R cycle (Retrieve, Reuse, Revise, Retain) with Jaccard similarity and BLAKE3-identified retained cases. Prolog received flat-term Robinson unification and bounded SLD resolution with positional \texttt{?N} variables, replacing the pattern-matching approximations common in toy implementations. These upgrades are not cosmetic: each is verified by tests that would fail against the simplified version.

\textbf{6. A taxonomy of eight Fortune-5 deployment patterns enabled by microsecond-latency certified reasoning.}
At 2--50 microseconds per breed execution on 2026 hardware, the cognition layer operates below the threshold of network latency, human perception, and most downstream decision pipelines. Chapter~6 identified eight deployment patterns---from edge inference at network ingress to regulatory audit trail generation---that become feasible when reasoning is both fast enough to inline and certified enough to cite. The BLAKE3 receipt chain provides the audit trail; the OCEL conformance gate provides the process certificate; the WASM compilation target provides the deployment portability.

\section{The Central Claim Restated}
\label{sec:central-claim}

The thesis has argued that classical AI systems were not wrong---they were unverifiable. The field did not abandon MYCIN because certainty factors were theoretically unsound; it abandoned MYCIN because there was no way to know, in production, whether MYCIN was actually applying Shortliffe and Buchanan's formula or a degraded approximation of it. The same applies to every architecture in the classical canon: the algorithms were specified, but their executions were opaque.

The \texttt{wasm4pm} framework makes them verifiable. Every execution is witnessed by an OCEL log that records the actual sequence of lifecycle phases traversed. That log is checked against a declared lifecycle model by a DFA conformance checker that refuses non-conforming executions. The entire witnessed execution is covered by a BLAKE3 receipt that makes the attestation tamper-evident. The receipt is not a log entry; it is a cryptographic commitment to process behavior.

This constitutes operational falsifiability in the Popperian sense. You cannot fake compliance with the Prolog grandparent test without implementing shared-variable unification. You cannot fake STRIPS plan soundness without simulating action application against state. You cannot fake MYCIN's CF threshold without evaluating the Shortliffe-Buchanan formula. The falsification suite does not test these properties through indirect proxies---it tests them through the exact mathematical conditions their inventors specified. A system that passes the suite either implements the original algorithm or has found a distinct algorithm that satisfies every one of its formal invariants, which would itself be a contribution.

\section{The Civilization-Scale Claim}
\label{sec:civilization}

This chapter would be incomplete without stating the larger claim directly, because the contributions enumerated above are not their own justification.

Certified reasoning at 2--50 microseconds, composition via shared contract, and BLAKE3 receipt chains as civil evidence are the components of \emph{reasoning infrastructure}. Not reasoning tools. Not reasoning applications. Infrastructure---in the same sense that the International System of Units is measurement infrastructure, or that TCP/IP is communication infrastructure. Infrastructure is what allows engineers to build on top of shared, trusted, auditable foundations without re-deriving the foundations each time.

The analogy is not rhetorical. The same way structural engineers certify bridges against load specifications, and pharmacologists certify drugs against bioequivalence standards, AI engineers can now certify reasoning chains against declared process models. The certificate is not a summary or a confidence score; it is a receipt whose hash coverage makes it unforgeable. The process model is not a diagram or a description; it is a DFA whose conformance checker rejects deviations at runtime.

What was missing from the classical AI era was not correctness---many of those systems were correct in controlled settings---but certifiability at scale. The OCEL provability layer provides the missing component. Whether that layer is used to certify medical decision support, financial audit trails, or regulatory compliance pipelines is a question of application, not architecture. The architecture is now available.

\section{Future Work}
\label{sec:future-work}

Six directions follow directly from the contributions of this dissertation.

\textbf{Cross-breed OCEL discovery at scale.} The current OCEL layer records individual breed executions. A natural extension is to accumulate OCEL logs across 10 million or more executions and apply process discovery algorithms---alpha miner, inductive miner, heuristic miner---to the resulting corpus. The goal is not to rediscover the declared lifecycle models but to detect \emph{variant distributions}: which lifecycle variants actually occur in production, how frequently, and under what input conditions. This is a direct application of the van der Aalst conformance framework to AI execution behavior.

\textbf{OCPQ-based conformance for complex lifecycle constraints.} The current DFA conformance checker handles linear phase sequences. The \texttt{ocpq} crate (object-centric process querying) supports path queries over object-centric event graphs, enabling constraints such as ``the Retain phase of every CBR execution must be preceded by a successful Revise phase on the same case object.'' Integrating \texttt{ocpq} with the OCEL layer would elevate conformance checking from phase ordering to cross-object causal consistency.

\textbf{Neural-symbolic integration with validated classical breeds as the reasoning layer.} Large language models exhibit brittleness on tasks requiring formal deduction, constraint satisfaction, and rule-consistent inference. The thirteen validated breeds provide certified reasoning modules that can serve as the symbolic layer in a hybrid architecture: LLM for natural language interface and abductive hypothesis generation, classical breed for certified deductive closure. The BLAKE3 receipt chain makes the boundary between the two layers auditable.

\textbf{Formal verification of the OCEL conformance checker itself.} The \texttt{validate\_ocel\_alignment} function is currently verified by 7 Rust tests and the broader adversarial suite. A stronger guarantee would come from a Coq or Lean proof that the DFA implementation correctly decides membership in the language defined by the lifecycle model. This is tractable: the lifecycle models are small finite automata, and the checker is a straightforward DFA simulation. The proof would elevate the conformance gate from ``tested'' to ``verified.''

\textbf{Regulatory framework mapping.} The FDA 21 CFR Part 11 requirements for electronic records, the EMA Annex 11 requirements for computerized systems, and the SEC Rule 17a-4 requirements for audit trails all specify tamper-evidence and traceability properties. The BLAKE3 receipt chain satisfies these properties structurally. Future work should produce a formal mapping from receipt chain fields to specific regulatory requirements, enabling direct citation of receipts in audit submissions.

\textbf{The \texttt{wpm cognition conformance} subcommand.} The \texttt{wpm} CLI currently supports breed execution and receipt inspection. A \texttt{conformance} subcommand that accepts an accumulated OCEL corpus and computes the four van der Aalst process quality dimensions---fitness, precision, generalization, and simplicity---over that corpus would close the loop between the process mining literature and the AI certification framework. This is the final component of the reasoning infrastructure: not just certifying individual executions, but auditing the aggregate process behavior of a deployed breed over time.

\section{Closing}
\label{sec:closing}

The history of artificial intelligence is a history of systems that could reason but could not prove it. MYCIN could recommend treatment; it could not prove its recommendation was derived from the declared certainty factor formula. STRIPS could find plans; it could not prove its plans were sound with respect to the operator specifications. Prolog could derive conclusions; it could not prove its derivation followed Robinson's unification algorithm rather than a local approximation of it.

The OCEL provability layer, the BLAKE3 receipt chain, and the 655-test falsification suite together constitute a proof that this gap can be closed. The gap was never algorithmic; the algorithms were sound. The gap was architectural: there was no mechanism to attest, in production, that the algorithm specified was the algorithm executed.

That mechanism now exists. It runs in WebAssembly, compiles to browser and server targets, executes in microseconds, and produces receipts whose hash coverage makes them unforgeable. The classical AI canon---seventy years of architectures, formalisms, and theorems---is now deployable as certified infrastructure.

For seventy years, we built machines that could reason. This work proves they do.
